\documentclass[12pt]{article}
\usepackage{amsmath,amssymb,amsthm,hyperref}
\usepackage{mathtools}
\usepackage[margin=1in]{geometry}

% Theorem environments
\newtheorem{theorem}{Theorem}[section]
\newtheorem{proposition}[theorem]{Proposition}
\newtheorem{lemma}[theorem]{Lemma}
\newtheorem{corollary}[theorem]{Corollary}
\newtheorem{claim}[theorem]{Claim}
\theoremstyle{definition}
\newtheorem{definition}[theorem]{Definition}
\newtheorem{remark}[theorem]{Remark}

% Macros
\newcommand{\N}{\mathbb{N}}
\newcommand{\Z}{\mathbb{Z}}
\newcommand{\R}{\mathbb{R}}
\newcommand{\Exp}{\mathbb{E}}
\newcommand{\SN}{\mathcal{S}_N}

\title{The Largest Sidon Subset of the Perfect Squares
$\{1^2, 2^2, \ldots, N^2\}$}
\author{}
\date{April 2026}

\begin{document}
\maketitle

\begin{abstract}
A \emph{Sidon set} (also called a $B_2$ set) is a set of integers in which
all pairwise sums are distinct.  Let $F(N)$ denote the maximum size of a
Sidon subset of $\mathcal{S}_N = \{1^2, 2^2, \ldots, N^2\}$.  We prove
\[
  \Omega\!\left(\frac{N^{2/3}}{(\log N)^{1/3}}\right)
  \;\le\;
  F(N)
  \;\le\;
  O\!\left(\frac{N}{(\log N)^{1/4}}\right).
\]
In particular $F(N) = \omega(N^{1/2})$.  The upper bound uses the
Landau--Ramanujan theorem on the density of sums of two squares.
The lower bound is a probabilistic alteration construction.
The question of whether $F(N) = N^{1-o(1)}$ is not settled by these bounds.
\end{abstract}

\tableofcontents

%--------------------------------------------------------------------
\section{Introduction and Statement of Results}
%--------------------------------------------------------------------

\begin{definition}
A set $A$ of integers is a \emph{Sidon set} (or \emph{$B_2$ set}) if
whenever $a,b,c,d \in A$ satisfy $a+b = c+d$ we have $\{a,b\} = \{c,d\}$
as multisets.
\end{definition}

Let $N \ge 1$ and define
\[
  \SN \;=\; \{k^2 : 1 \le k \le N\}, \quad |\SN| = N.
\]
Set
\[
  F(N) \;:=\; \max\bigl\{|A| : A \subseteq \SN,\; A \text{ is Sidon}\bigr\}.
\]

\begin{theorem}[Main Theorem]\label{thm:main}
\[
  \Omega\!\left(\frac{N^{2/3}}{(\log N)^{1/3}}\right)
  \;\le\;
  F(N)
  \;\le\;
  O\!\left(\frac{N}{(\log N)^{1/4}}\right).
\]
\end{theorem}

\begin{corollary}\label{cor:not-N}
$F(N) = \omega(N^{1/2})$.
\end{corollary}

\begin{remark}\label{rem:N1-o1}
Note that $N/(\log N)^{1/4} = N^{1-o(1)}$ (since the exponent
$1 - \frac{\log\log N}{4\log N} \to 1$), so the upper bound
$O(N/(\log N)^{1/4})$ is entirely consistent with $F(N) = N^{1-o(1)}$.
The question of whether $F(N) = N^{1-o(1)}$ is \emph{not settled}
by Theorem~\ref{thm:main}.
\end{remark}

\subsection{Historical context}

The classical Sidon problem asks for the maximum size of a Sidon subset
of $\{1,\ldots,M\}$; the answer is $(1+o(1))M^{1/2}$ (Erd\H{o}s--Tur\'an
upper bound \cite{ET41}, Singer difference set lower bound \cite{Sin38}).
Our host set $\SN$ has $N$ elements and lies in $\{1,\ldots,N^2\}$, and the
Sidon condition on sums $a_i^2+a_j^2$ interacts with the analytic structure
of the two-squares representation function~$r_2$.

%--------------------------------------------------------------------
\section{Preliminaries}
%--------------------------------------------------------------------

\begin{definition}
For $n \ge 0$ define
\[
  r_2(n) = \#\{(a,b)\in\Z^2 : a^2+b^2=n\}, \quad
  r_2^+(n;N) = \#\{(a,b)\in\mathbb{N}^2 : a^2+b^2=n,\,1\le a\le b\le N\}.
\]
\end{definition}

\begin{lemma}[Second moment of $r_2$; {\cite[Ch.~18]{HW08}}]\label{lem:r2-second}
As $M\to\infty$,
\[
  \sum_{n=0}^{M} r_2(n)^2 \;\sim\; 4\pi M \log M.
\]
\end{lemma}

\begin{proof}
The sum equals the number of 4-tuples $(a,b,c,d)\in\Z^4$ with
$a^2+b^2=c^2+d^2\le M$.  Via the factorisation $r_2(n)=4(d_1(n)-d_3(n))$
(where $d_i$ counts divisors $\equiv i\pmod{4}$) and standard Dirichlet
series methods, one obtains the stated asymptotic; see \cite[Theorem~339]{HW08}.
\end{proof}

\begin{lemma}[Landau--Ramanujan; {\cite{Land08}}]\label{lem:LR}
As $M\to\infty$,
\[
  \#\{n \le M : r_2(n) > 0\} \;\sim\; K \cdot \frac{M}{\sqrt{\log M}},
\]
where $K = \frac{1}{\sqrt{2}}\prod_{p \equiv 3(4)}(1-p^{-2})^{-1/2}
\approx 0.7642\ldots$ is the Landau--Ramanujan constant.
\end{lemma}

An integer $n$ is a sum of two squares if and only if every prime
$p \equiv 3\pmod{4}$ divides $n$ to an even power; see \cite[Thm.~366]{HW08}.
Lemma~\ref{lem:LR} follows from this characterisation via a standard Euler
product computation and a Tauberian theorem.

%--------------------------------------------------------------------
\section{Upper Bound}\label{sec:upper}
%--------------------------------------------------------------------

\begin{theorem}\label{thm:upper}
$F(N) \le O\!\left(N/(\log N)^{1/4}\right)$.
\end{theorem}

\begin{proof}
Let $A = \{a_1^2,\ldots,a_k^2\} \subseteq \SN$ be a Sidon set with
$1 \le a_1 < \cdots < a_k \le N$ and $k = |A|$.

\medskip
\noindent\textbf{Step 1.}
Since $A$ is Sidon, the $\binom{k+1}{2}$ values
\[
  \Sigma(A) = \bigl\{a_i^2+a_j^2 : 1 \le i \le j \le k\bigr\}
\]
are pairwise distinct.  Each element of $\Sigma(A)$ is a sum of two
positive squares bounded by $N$, so
\[
  \Sigma(A) \;\subseteq\; R_N := \{n \le 2N^2 : r_2(n) > 0\}.
\]

\noindent\textbf{Step 2.}
By Lemma~\ref{lem:LR} with $M = 2N^2$,
\[
  |R_N| \;\le\; K \cdot \frac{2N^2}{\sqrt{\log(2N^2)}} \cdot (1+o(1))
  \;=\; \frac{2K\,N^2}{\sqrt{2\log N}} \cdot (1+o(1))
  \;=\; \frac{K\sqrt{2}\,N^2}{\sqrt{\log N}} \cdot (1+o(1)).
\]

\noindent\textbf{Step 3.}
Since all $\binom{k+1}{2}$ elements of $\Sigma(A)$ lie in $R_N$,
\[
  \frac{k(k+1)}{2} \;\le\; |R_N| \;\le\; C\,\frac{N^2}{(\log N)^{1/2}}
\]
for an absolute constant $C > 0$.  Hence $k^2 \le 2C\,N^2/(\log N)^{1/2}$,
giving $k \le \sqrt{2C}\,N/(\log N)^{1/4}$.
\end{proof}

%--------------------------------------------------------------------
\section{Lower Bound}\label{sec:lower}
%--------------------------------------------------------------------

\begin{theorem}\label{thm:lower}
$F(N) \ge \Omega\!\left(N^{2/3}/(\log N)^{1/3}\right)$.
\end{theorem}

We use the probabilistic method with alteration.

\begin{definition}
A \emph{collision} is an ordered quadruple $(a,b,c,d) \in \{1,\ldots,N\}^4$
with $a < b$, $c < d$, $\{a,b\} \ne \{c,d\}$, and $a^2+b^2 = c^2+d^2$.
Let $\mathcal{C}$ be the set of all collisions.
\end{definition}

\begin{lemma}\label{lem:C-bound}
$|\mathcal{C}| \le 8\pi N^2\log N \cdot (1+o(1))$.
\end{lemma}

\begin{proof}
For each $n \in \{2,\ldots,2N^2\}$, the collisions at level $n$ number
$\binom{r_2^+(n;N)}{2} \le \frac{1}{2}r_2^+(n;N)^2 \le \frac{1}{2}r_2(n)^2$.
Summing and applying Lemma~\ref{lem:r2-second} with $M = 2N^2$:
\[
  |\mathcal{C}| \;\le\; \frac{1}{2}\sum_{n=0}^{2N^2} r_2(n)^2
  \;\sim\; \frac{1}{2}\cdot 4\pi \cdot 2N^2 \cdot \log(2N^2)
  \;\sim\; 8\pi N^2\log N. \qedhere
\]
\end{proof}

\begin{proof}[Proof of Theorem~\ref{thm:lower}]
Let $\lambda > 0$ be a parameter to be chosen.  Include each
$a \in \{1,\ldots,N\}$ independently in a random set $\mathcal{A}$
with probability $\lambda$.  Let $X = |\mathcal{A}|$ and let $Y$ be
the number of collisions in $\mathcal{A}$.

\noindent\textbf{Expectations:}
\[
  \Exp[X] = N\lambda,
  \qquad
  \Exp[Y] = |\mathcal{C}|\lambda^4 \le 8\pi N^2\log N\cdot\lambda^4\cdot(1+o(1)).
\]

\noindent\textbf{Alteration:}  For each collision in $\mathcal{C}$ whose four
elements all lie in $\mathcal{A}$, mark one of the four for deletion; then
delete the union of all marked elements from $\mathcal{A}$, obtaining a
surviving index set~$\mathcal{A}^*$.  Because every collision is hit
by this procedure, the square set $\{a^2 : a \in \mathcal{A}^*\}$ is a
Sidon subset of $\SN$.  At most one element is deleted per collision,
so $|\mathcal{A}^*| \ge X - Y$.  Hence
\[
  \Exp[|\mathcal{A}^*|] \;\ge\;
  N\lambda - 8\pi N^2\log N\cdot\lambda^4\cdot(1+o(1)).
\]

\noindent\textbf{Optimising $\lambda$:}
Maximise $g(\lambda) = N\lambda - 8\pi N^2\log N\cdot\lambda^4$.
Setting $g'(\lambda) = N - 32\pi N^2\log N\cdot\lambda^3 = 0$ gives
\[
  \lambda^* \;=\; \left(\frac{1}{32\pi N\log N}\right)^{1/3}
  \;=\; \frac{1}{(32\pi)^{1/3}\,N^{1/3}(\log N)^{1/3}}.
\]
Since $8\pi N\log N\cdot(\lambda^*)^3
= 8\pi N\log N\cdot\frac{1}{32\pi N\log N} = \tfrac{1}{4}$, we get
\[
  g(\lambda^*) = N\lambda^*\!\left(1 - 8\pi N\log N\cdot(\lambda^*)^3\right)
  = N\lambda^*\!\left(1 - \tfrac{1}{4}\right)
  = \frac{3}{4}\,N\lambda^*
  = \frac{3}{4(32\pi)^{1/3}}\cdot\frac{N^{2/3}}{(\log N)^{1/3}}.
\]

\noindent\textbf{Conclusion:}
With $\lambda = \lambda^*$,
\[
  \Exp[|\mathcal{A}^*|] \;\ge\;
  \frac{3}{4(32\pi)^{1/3}}\cdot\frac{N^{2/3}}{(\log N)^{1/3}}\cdot(1+o(1)).
\]
Since the expectation is an average, there exists a realisation with
\[
  |\mathcal{A}^*| \;\ge\; \Exp[|\mathcal{A}^*|]
  \;\ge\; \frac{3}{4(32\pi)^{1/3}}\cdot\frac{N^{2/3}}{(\log N)^{1/3}}\cdot(1+o(1)),
\]
giving $F(N) \ge \Omega(N^{2/3}/(\log N)^{1/3})$.
\end{proof}

%--------------------------------------------------------------------
\section{Proof of the Main Results}
%--------------------------------------------------------------------

\begin{proof}[Proof of Theorem~\ref{thm:main}]
The lower bound is Theorem~\ref{thm:lower}.
The upper bound is Theorem~\ref{thm:upper}.
\end{proof}

\begin{proof}[Proof of Corollary~\ref{cor:not-N}]
Theorem~\ref{thm:lower} gives $F(N) \ge \Omega(N^{2/3}/(\log N)^{1/3})$.
Since $N^{2/3}/(\log N)^{1/3} \gg N^{1/2}$ (their ratio $N^{1/6}/(\log N)^{1/3} \to \infty$),
this is $\omega(N^{1/2})$.
\end{proof}

%--------------------------------------------------------------------
\section{Discussion and Open Problems}
%--------------------------------------------------------------------

\subsection{Summary}

We have proved:
\begin{enumerate}
  \item $F(N) = \omega(N^{1/2})$: the maximum Sidon subset of $\SN$
        grows strictly faster than $N^{1/2}$.
  \item $F(N) \le O(N/(\log N)^{1/4})$: it grows no faster than
        $N$ up to a logarithmic factor.
\end{enumerate}

\noindent\textbf{The question of whether $F(N) = N^{1-o(1)}$ is not resolved.}
The upper bound $O(N/(\log N)^{1/4})$ is itself $N^{1-o(1)}$
(since $(\log N)^{1/4} = N^{o(1)}$, so the exponent of $N$ approaches $1$),
and thus does not rule out $F(N) = N^{1-o(1)}$.
The lower bound $\Omega(N^{2/3}/(\log N)^{1/3})$ is $N^{2/3+o(1)}$,
which is consistent with the true answer being either $N^{1-o(1)}$ or
strictly smaller.  Closing this gap is the central open problem.

\subsection{Sharpness}

\begin{itemize}
\item The Landau--Ramanujan bound in Step~2 of the upper bound proof is sharp:
  the number of sums-of-two-squares up to $M$ is indeed $\Theta(M/(\log M)^{1/2})$.
\item The alteration method is flexible but not tight; a construction via
  Sidon sets in $\mathbb{F}_p$ or difference sets could potentially give a
  better explicit lower bound.
\end{itemize}

\subsection{Why the $N^{1-o(1)}$ question remains open}

Neither bound in Theorem~\ref{thm:main} is strong enough to determine
whether $\alpha = 1$.

\paragraph{The upper bound does not exclude $\alpha = 1$.}
The proof of Theorem~\ref{thm:upper} uses only the inequality
$\tbinom{|A|+1}{2} \le |R_N|$, where $|R_N| \sim KN^2/\sqrt{\log N}$.
This gives $F(N) \le O(N/(\log N)^{1/4})$, but
\[
  \frac{\log\bigl(N/(\log N)^{1/4}\bigr)}{\log N}
  \;=\; 1 - \frac{\log\log N}{4\log N} \;\longrightarrow\; 1,
\]
so the bound is $N^{1-o(1)}$ and places no constraint on
the true exponent $\alpha$ being strictly less than~$1$.
No variant of this counting argument can yield a power-saving bound
$O(N^{1-\varepsilon})$ for any fixed $\varepsilon > 0$, because the
argument uses only the \emph{cardinality} of $R_N$ and not the deeper
additive structure of sums of two squares.

\paragraph{The lower bound does not reach $N^{1-o(1)}$.}
The probabilistic construction gives $F(N) \ge \Omega(N^{2/3}/(\log N)^{1/3})$,
an exponent of $2/3 + o(1)$.  A construction of size $N^{1-o(1)}$
would require a Sidon subset of $\SN$ of size $N/(\log N)^{o(1)}$, which
is far beyond the reach of current methods.

\medskip\noindent
Resolving whether $\alpha = 1$ therefore requires either:
\begin{enumerate}
  \item a new upper bound exploiting the multiplicative structure of the
        integers beyond a sums-of-two-squares count, sufficient to give
        $F(N) = O(N^{1-\varepsilon})$ for some fixed $\varepsilon > 0$; or
  \item an algebraic or number-theoretic construction of a Sidon subset of
        $\SN$ of size $N/(\log N)^{o(1)}$.
\end{enumerate}
Both appear to require genuinely new ideas.

\subsection{Open problems}

\begin{itemize}
\item \textbf{Determine the exponent $\alpha$} (assuming $F(N) = N^{\alpha+o(1)}$).
  The current bounds show $\alpha \in [2/3,1]$ but do not pin down the
  value, or even whether $\alpha < 1$.
\item \textbf{Improve the upper bound.}  Prove $F(N) = O(N^{1-\varepsilon})$
  for some fixed $\varepsilon > 0$, or else exhibit a Sidon subset of $\SN$
  of size $\omega(N^{2/3})$.
\item \textbf{Explicit construction.}  Find a deterministic construction of
  a Sidon subset of $\SN$ of size $\Omega(N^{2/3})$.
\end{itemize}

%--------------------------------------------------------------------
\begin{thebibliography}{99}

\bibitem{ET41}
P.~Erd\H{o}s and P.~Tur\'an,
\textit{On a problem of Sidon in additive number theory, and on some related problems},
J. London Math. Soc. \textbf{16} (1941), 212--215.

\bibitem{HW08}
G.~H. Hardy and E.~M. Wright,
\textit{An Introduction to the Theory of Numbers},
6th ed., Oxford University Press, Oxford, 2008.

\bibitem{Land08}
E.~Landau,
\textit{\"Uber die Einteilung der positiven ganzen Zahlen in vier Klassen
nach der Mindestzahl der zu ihrer additiven Zusammensetzung erforderlichen Quadrate},
Arch. Math. Phys. \textbf{13} (1908), 305--312.

\bibitem{Sin38}
J.~Singer,
\textit{A theorem in finite projective geometry and some applications to number theory},
Trans. Amer. Math. Soc. \textbf{43} (1938), 377--385.

\bibitem{Ten}
G.~Tenenbaum,
\textit{Introduction to Analytic and Probabilistic Number Theory},
3rd ed., American Mathematical Society, Providence, RI, 2015.

\end{thebibliography}

\end{document}
